% !TeX root=../main.tex
% در این فایل، عنوان پایان‌نامه، مشخصات خود، متن تقدیمی‌، ستایش، سپاس‌گزاری و چکیده پایان‌نامه را به فارسی، وارد کنید.
%%%%%%%%%%%%%%%%%%%%%%%%%%%%%%%%%%%%
% --- نکتهٔ نویسنده: موارد زیر (دانشگاه، دانشکده، گروه، اساتید، داوران، نام و شمارهٔ دانشجویی) با مقادیر پیش‌فرض/جای‌گذار پر شده‌اند و باید پیش از ارائهٔ نهایی، توسط شما با اطلاعات واقعی جایگزین شوند. ---
\university{دانشگاه تهران}
\college{پردیس دانشکده‌های فنی}
\faculty{دانشکدهٔ مهندسی برق و کامپیوتر}
\department{گروه هوش مصنوعی}
\subject{مهندسی کامپیوتر}
\field{هوش مصنوعی}
% عنوان پایان‌نامه
\title{ریزتنظیم مدل‌های زبانی برای ناوردایی معنایی در فضای نهان خودرمزگذارهای پراکنده}
% نام استاد(ان) راهنما -- جای‌گذار
\firstsupervisor{دکتر مصطفی توسلی‌پور}
\firstsupervisorrank{استادیار}
%\secondsupervisor{دکتر راهنمای دوم}
%\secondsupervisorrank{استادیار}
% استاد مشاور -- در صورت نداشتن مشاور، این دو خط را غیرفعال کنید
%\firstadvisor{[نام استاد مشاور]}
%\firstadvisorrank{استادیار}
% داوران -- جای‌گذار
\internaljudge{دکتر یداله یعقوب‌زاده}
\internaljudgerank{استادیار}
% \externaljudge{[نام داور خارجی]}
% \externaljudgerank{استادیار}
% \externaljudgeuniversity{[دانشگاه داور خارجی]}
% نمایندهٔ تحصیلات تکمیلی -- جای‌گذار
% \graduatedeputy{[نام نماینده]}
% \graduatedeputyrank{دانشیار}
% نام و شمارهٔ دانشجویی -- جای‌گذار
\name{رضا}
\surname{چهرقانی}
\studentID{۸۱۰۱۰۱۴۰۱}
\thesisdate{۱۴۰۵}
%%%%%%%%%%%%%%%%%%%%%%%%%%%%%%%%%%%%%%%%%%%%%%%%%%%%
\dedication
{
{\Large تقدیم به:}\\
\begin{flushleft}{
	\huge
	هرکس که در طلب دانش، شکیبایی را نیز آموخت
}
\end{flushleft}
}
%% متن قدردانی -- جای‌گذار؛ لطفاً با نام واقعی استاد راهنما و دیگران تکمیل کنید.
\acknowledgement{
از استاد راهنمای گران‌قدرم، دکتر توسلی‌پور، برای راهنمایی‌های ارزنده و صبورانه در سرتاسر این پژوهش، صمیمانه سپاس‌گزارم.

از داوران محترم این پایان‌نامه، برای وقتی که صرف مطالعه و ارزیابی این نوشتار کردند، کمال امتنان را دارم.

و سپاس ویژه از خانواده‌ام، برای حمایت بی‌دریغشان در طول این مسیر.
}
%%%%%%%%%%%%%%%%%%%%%%%%%%%%%%%%%%%%
\fa-abstract{
مدل‌های زبانی بزرگ، بازنمایی‌های برداری خود را عمدتاً با هدف پیش‌بینی نشانهٔ بعدی می‌آموزند، هدفی که لزوماً بازنمایی جمله را در راستای سنجه‌های معنایی مورد نظر انسان شکل نمی‌دهد. مجموعه‌دادهٔ
\lr{PAWS}
این شکاف را آشکارا نشان می‌دهد: جفت‌جمله‌هایی با همپوشانی واژگانی بسیار بالا که برخی هم‌معنی و برخی ناهم‌معنی‌اند.% بررسی مقدماتی این پایان‌نامه نشان داد در یک مدل زبانی پیش‌آموزش‌دیده، فاصلهٔ بازنمایی این دو دسته جفت، به‌قدر کافی از هم متمایز نیست.

این پایان‌نامه روشی را برای ریزتنظیم کارآمد پارامتر یک مدل زبانی پیشنهاد می‌دهد که در آن، به‌جای اعمال مستقیم یک تابع هزینهٔ ناوردایی (که جفت‌های هم‌معنی را نزدیک و جفت‌های ناهم‌معنی را دور از هم نگاه می‌دارد) بر فضای خام جریان باقیماندهٔ مدل، این تابع هزینه بر کدهای نهان حاصل از عبور فعال‌سازی‌ها از یک خودرمزگذار پراکندهٔ از پیش‌آموزش‌دیده و ثابت اعمال می‌شود. انگیزهٔ این انتخاب، پدیدهٔ برهم‌نهی ویژگی‌ها در فضای خام مدل‌های زبانی است که اعمال مستقیم چنین هدفی را در آن فضا غیرقابل‌پیش‌بینی می‌کند.% برای مهار ریسک رانش خودرمزگذار در حین آموزش، یک منظم‌ساز صریح--مبتنی بر خطای بازسازی خود خودرمزگذار--به تابع هزینهٔ کلی افزوده شده است.

این روش بر روی مدل
\lr{Pythia-70m-deduped}
با تطبیق‌دهندهٔ لورا پیاده‌سازی و در یک شبکهٔ آزمایشی سیستماتیک--شامل چهار شرط آزمایشی (کنترل، ناوردایی در فضای خام، ناوردایی در فضای خودرمزگذار، و نسخهٔ تکمیل‌شدهٔ آن)، پنج مقدار برای وزن نسبی تابع هزینهٔ ناوردایی، و ده حالت متفاوت برای انتخاب لایهٔ هدف--با مجموع ۱۶۰ اجرای مستقل، ارزیابی شده است.

نتایج نشان می‌دهند اعمال تابع هزینهٔ ناوردایی در فضای خودرمزگذار، به‌طور پایدار، جداسازی بهتری میان جفت‌های هم‌معنی و ناهم‌معنی (بر حسب سطح زیر منحنی مشخصهٔ عملکرد گیرنده) نسبت به اعمال همان تابع هزینه در فضای خام یا نبود هرگونه هدف ناوردایی تولید می‌کند. این بهبود به‌شدت به عمق لایهٔ هدف (لایه‌های میانی-به-بعد، نه لایه‌های ابتدایی) و به وزن تابع هزینه وابسته است، و اعمال هم‌زمان چند لایه اثری فراجمع‌شونده دارد. بااین‌همه، بالاترین سطوح این بهبود با ناپایداری گذرا در آموزش و با افت اعتبار بازسازی خودرمزگذار همراه‌اند؛ منظم‌ساز پیشنهادی این هزینه را در بازهٔ وزن‌های پایین به‌طور کامل، و در بازهٔ وزن‌های بالا تنها تا حدی مهار می‌کند. بر همین اساس، این پایان‌نامه یک پیکربندی متعادل--با وزن میانه‌رو‌تر--را معرفی می‌کند که تقریباً به همان بهبود دست می‌یابد، بدون تحمیل این هزینه. این پایان‌نامه با ترسیم این تصویر متعادل--به‌جای گزارش صرف بهترین نتیجهٔ خام--کاربرد تازه‌ای برای خودرمزگذارهای پراکنده، فراتر از نقش متداول آن‌ها در تفسیرپذیری پس‌ازآموزش، معرفی می‌کند.
}
\keywords{مدل‌های زبانی، خودرمزگذار پراکنده، ریزتنظیم کارآمد پارامتر، یادگیری تقابلی، تشخیص بازنویسی، تفسیرپذیری مکانیستی}
%%%%%%%%%%%%%%%%%%%%%%%%%%%%%%%%%%%%%%%%%%%%%%%%%%%%%%
